\documentclass[11pt,a4paper]{article}
\usepackage[utf8]{inputenc}
\usepackage[T1]{fontenc}
\usepackage{amsmath,amssymb,amsfonts}
\usepackage{hyperref}
\usepackage{cite}
\usepackage{listings}

\title{Matrix State Resolution v2.2: \\ Formally Verified and Zero-Knowledge Accelerated Consensus}
\author{Shane Jaroch \\ \texttt{chown\_tee@proton.me}}
\date{June 2026}

\begin{document}

\maketitle

\begin{abstract}
Matrix State Resolution v2 is the bedrock of decentralized consensus in the Matrix protocol. However, as rooms scale to hundreds of thousands of events and deep federation forks, the algorithm's performance degrades and edge-case vulnerabilities emerge. We present State Resolution v2.2, which introduces the ``Transitive'' supplemental merge to eliminate the Power Level Replay vulnerability (CVE-2026-XXXXX) while maintaining Spec-compliant behavior. Furthermore, we describe a high-performance implementation utilizing Roaring Bitmaps and memoized auth contexts, and a Lean 4 formalization that enables Zero-Knowledge STARK proofs of room state.
\end{abstract}

\section{Introduction}
State Resolution is the process by which multiple homeservers in a Matrix room reach a common view of the room state despite participating in a distributed, asynchronous Directed Acyclic Graph (DAG).

\section{The Power Level Replay Vulnerability}
In State Resolution v2.1 (MSC4297), the introduction of a supplemental merge to stabilize room state inadvertently opened a window for ``Time-Travel'' attacks. A demoted administrator could craft a malicious event citing their historical (elevated) power levels.

\section{State Resolution v2.2: The Transitive Fix}
State Res v2.2 solves this by restricting the supplemental merge to \textit{ancestor} power levels only. An event $e$ can only authenticate against a resolved power level $PL_{resolved}$ if $PL_{resolved}$ is already present in the recursive auth chain of $e$.

\section{Performance Optimizations}
\subsection{Memoized Auth Contexts}
Traditional implementations perform a full Graph Walk (BFS) for every event in the conflict set. In v2.2, we memoize the set of auth-relevant events (Create, Membership, Join Rules, Power Levels) as we traverse the topological sort.

\subsection{Roaring Bitmaps for Reachability}
By mapping Event IDs to integer indices, we can utilize SIMD-accelerated bitwise operations (Roaring Bitmaps) to calculate auth chain differences and reachability in $O(1)$ or $O(\log N)$ time.

\section{Formal Verification and ZK-STARKs}
The Ruma-Lean implementation is formally verified using the Lean 4 theorem prover. This verification enables the generation of STARK proofs, allowing a server to prove the validity of a room's state to a light client without the client needing to re-execute the entire DAG.

\end{document}
